% report/sections/04_task2_results.tex -- Task 2 results (C3, Member 2)
% Every figure and table comes from `python -m project4.experiments.task2_cost_sweep`
% (files in project4/figures and project4/results). Do not type numbers by hand.
% Discussion sentences that depend on the measured numbers are marked \todo{...}:
% check each against the real run, then rewrite or delete the marker.
\providecommand{\todo}[1]{\textbf{[TODO: #1]}}
\providecommand{\cFP}{c_{\mathrm{FP}}}
\providecommand{\cFN}{c_{\mathrm{FN}}}
\providecommand{\cI}{c_{\mathrm{I}}}

\section{Task 2 results: accuracy-based versus loss-minimising decisions}
\label{sec:task2-results}

\subsection{Experimental setup}
\paragraph{Data and splits.}
We use UCI Spambase (4601 emails, 57 features, 1813 spam) and a synthetic two-class Gaussian data set with
$d=5$, $n=4601$, $\pi_1=0.39$, $\mu_0=0$, $\mu_1=0.9\cdot\mathbf 1$, $\Sigma_0=I$ and
$\Sigma_1=\tfrac12 I+\tfrac12\mathbf 1\mathbf 1^{\top}$. The covariances differ, so the true posterior is quadratic
and the exact Bayes risk can be computed. Each data set is split 60/20/20 into training, validation and test parts,
stratified by class. All results are means with standard errors over 20 repeated splits (seeds $0$--$19$).
Features are standardised with training statistics only. The test part is never used to choose anything.

\paragraph{Models.}
We fit logistic regression (Newton/IRLS with an L2 penalty on the weights, Section~\ref{sec:models}) and the
Gaussian generative classifier with shared covariance and with class-specific covariances
(Section~\ref{sec:t2}). The L2 strength $\lambda_2\in\{10^{-2},10^{-1.5},\dots,10^{3}\}$ and the Gaussian
shrinkage strength $\lambda\in\{10^{-3},10^{-2},0.03,0.1,0.3,1\}$ of Section~\ref{sec:t2-reg} are chosen on the validation
part by minimum log-loss. On the synthetic data we also report an \emph{oracle} that applies the same decision
rules to the exact posterior.

\paragraph{Costs.}
The loss matrix is that of Task~1: correct decisions cost $0$, a real email sent to spam costs $\cFP$, spam
delivered costs $\cFN$ and an inspection costs $\cI$, with the inspector always finding the true label. We sweep
the cost ratio $r=\cFP/\cFN$ on a log grid from $1/20$ to $20$ with
$\cFP=\sqrt r$ and $\cFN=1/\sqrt r$, and keep $\cI=0.1$ fixed. For this scheme the critical inspection cost
$\cFP\cFN/(\cFP+\cFN)=\sqrt r/(1+r)$ is at least $0.21$ on the whole grid, so $\cI<\cFP\cFN/(\cFP+\cFN)$ holds at
every ratio and the inspect band is never empty.

\paragraph{Decision rules.}
For each fitted posterior $p=P(y=1\mid x)$ we compare three rules.
\begin{description}[leftmargin=1.5em,labelsep=0.5em,itemsep=2pt]
  \item[Accuracy rule] send to spam iff $p\ge\tfrac12$, no inspection. This is the Bayes rule for 0--1 loss.
  \item[Bayes, two actions] send to spam iff $p\ge\cFP/(\cFP+\cFN)$, no inspection.
  \item[Bayes with inspection] deliver if $p\le\theta_0=\cI/\cFN$, send to spam if $p\ge\theta_1=1-\cI/\cFP$,
        inspect otherwise.
\end{description}
The reported quantity is the average realized loss per test email, $\frac1N\sum_n L[a_n,y_n]$, so an inspection
contributes $\cI$ whatever the true label.

\subsection{Average realized loss against the cost ratio}
Figure~\ref{fig:t2-loss} shows the loss of the three rules as the ratio varies, and
Tables~\ref{tab:t2-loss-spambase} and~\ref{tab:t2-loss-synthetic} give the values at three representative
ratios, $r\in\{0.1,1,10\}$.

\begin{figure}[ht]
  \centering
  \includegraphics[width=\textwidth]{task2_loss_spambase}\\[4pt]
  \includegraphics[width=\textwidth]{task2_loss_synthetic}
  \caption{Average realized test loss against $\cFP/\cFN$ (log scale), mean $\pm$ standard error over 20 splits,
  $\cI=0.1$. Top: Spambase. Bottom: synthetic data, where the dashed line is the oracle with inspection.
  Dotted vertical lines mark the ratios of the tables.}
  \label{fig:t2-loss}
\end{figure}

\begin{table}[ht]
  \centering\small
  \input{\resultsdir/task2_loss_spambase.tex}
  \caption{Spambase: average realized test loss, mean $\pm$ standard error over 20 splits.}
  \label{tab:t2-loss-spambase}
\end{table}

\begin{table}[ht]
  \centering\small
  \input{\resultsdir/task2_loss_synthetic.tex}
  \caption{Synthetic data: average realized test loss, mean $\pm$ standard error over 20 splits.}
  \label{tab:t2-loss-synthetic}
\end{table}

\subsection{Confusion matrices}
Tables~\ref{tab:t2-conf-spambase} and~\ref{tab:t2-conf-synthetic} give the mean confusion matrices of logistic
regression, with the inspect column, at the three representative ratios (counts per test split,
about 920 emails). Figure~\ref{fig:t2-reject} shows how the share of inspected emails changes with the ratio.
The matrices for the Gaussian models are produced by the same script and are listed in the appendix.

\begin{table}[ht]
  \centering\small
  \input{\resultsdir/task2_confusion_spambase_logistic.tex}
  \caption{Spambase, logistic regression: mean confusion matrices (true label $\times$ action).}
  \label{tab:t2-conf-spambase}
\end{table}

\begin{table}[ht]
  \centering\small
  \input{\resultsdir/task2_confusion_synthetic_logistic.tex}
  \caption{Synthetic data, logistic regression: mean confusion matrices (true label $\times$ action).}
  \label{tab:t2-conf-synthetic}
\end{table}

\begin{figure}[ht]
  \centering
  \includegraphics[width=\textwidth]{task2_reject_spambase}\\[4pt]
  \includegraphics[width=\textwidth]{task2_reject_synthetic}
  \caption{Share of emails held for inspection under the Bayes rule with inspection, against $\cFP/\cFN$.
  Top: Spambase. Bottom: synthetic data.}
  \label{fig:t2-reject}
\end{figure}

\subsection{Discussion}
\paragraph{The accuracy rule is optimal only at equal costs.}
The accuracy rule uses threshold $\tfrac12$, which equals the Bayes two-action threshold
$\cFP/(\cFP+\cFN)$ only when $\cFP=\cFN$. At $r=1$ the two rules therefore make identical decisions and
have identical loss, as the figures and tables show. Away from $r=1$ the accuracy rule is
suboptimal in expectation whenever the posterior is calibrated, and its loss rises as $|\log r|$ grows,
because it keeps treating both kinds of error alike while one of them becomes much more expensive.
\todo{state the measured loss gap between the accuracy rule and the Bayes two-action rule at $r=0.1$ and $r=10$
for the three models, on both data sets.}

\paragraph{Allowing inspection lowers the loss.}
Since $\cI$ is below the critical value at every ratio, the Bayes rule with inspection has a nonempty band and
a lower expected loss than any two-action rule, provided the posterior is calibrated. In the tables it should
therefore have the lowest loss in every column.
\todo{confirm this for every model and ratio, and report the reject rate at $r=0.1,1,10$ from
Figure~\ref{fig:t2-reject}. Say which kind of email is inspected: the confusion tables show how the inspect
column splits between real and spam emails.}

\paragraph{Comparison of the models.}
On the synthetic data the true covariances differ, so the class-specific Gaussian model is correctly specified
and should come closest to the oracle, while the shared-covariance Gaussian and logistic regression are
linear and cannot represent the quadratic boundary.
\todo{report the gap of each model to the oracle at the three ratios. On Spambase say which model has the lowest
loss and whether the class-specific Gaussian suffers from its $3306$ covariance parameters, using the chosen
shrinkage strengths in \texttt{results/task2\_hyperparams.csv}.}

\paragraph{Calibration matters.}
The thresholds $\theta_0,\theta_1$ are the Bayes thresholds only if the model's probabilities are calibrated
(Section~\ref{sec:t2}). A model with the same accuracy but saturated probabilities, as the Gaussian models can
produce on Spambase, can be well off at the optimal thresholds, which is the point of the calibration analysis
that follows. \todo{compare the logistic and Gaussian reject rates and losses and relate them to the
reliability plots of the calibration section, if the numbers support it.}

\paragraph{Limitations.}
The inspector is assumed perfect and $\cI$ is held fixed; the sensitivity to $\cI$ is the subject of
Task~5. Each test part has about 920 emails, so differences smaller than the standard errors in the tables
should not be interpreted.
